% mcm-paper source binding (non-rendering metadata)
% Repository HEAD: 66f49838fe3060c9c0286cb5aa87613ac20b9fc3
% Abstract content version: V0.4
% Recommended title: 海上人员运输的异构直升机路径与机队排班优化
% NOTE: formal filename remains 00_abstract.tex.

\begin{abstract}
海上设施人员运输同时受到异构直升机、座位与燃油、多机场闭环、海上着陆次数、人员取送、时间窗及多日机队排班等约束。针对题目三问，本文建立“统一架次可行性核—候选路线/时间列生成—整数主问题—词典序优化—独立验证”的递进式优化框架，在统一物理规则下完成单向出海、联合取送与有限机队多日调度。

针对问题一，构建有限物理路线池—人员流整数规划模型，将路线执行次数与请求组承运人数分离。候选路线由地理近邻设施簇、访问顺序枚举及燃油标号搜索生成，并结合解引导扩池与CP-SAT五级词典序优化完成求解。最终1600名人员全部完成运输，飞机使用时间为14759 min，共执行88架次，燃油消耗118693 kg。

针对问题二，构建物理路线池—服务选项动态人员流联合取送整数规划模型，以服务选项刻画不同人员在候选路线中的取送方式，并利用半开区间容量约束描述座位即时复用。通过确定性成对插入和解引导路线池扩展生成联合运输候选路线，再利用CP-SAT依次优化五级目标。最终4000名人员全部完成运输，飞机使用时间18963 min，共执行111架次，距离加权座位利用率84.95\%；相较独立OD运输基准，飞机使用时间、燃油消耗和架次数分别下降66.4\%、68.0\%和69.0\%，座位利用率提高57.5个百分点。

针对问题三，构建时间可行架次列—具体飞机指派—服务选项多日机队排班整数规划模型，将实际起飞时刻、任务时间窗、30 min周转及单机时序冲突纳入统一约束。根据任务时间窗和路线持续时间生成临界起飞时刻及时间可行架次列，并采用A--C三阶段词典序求解：先确定非临时任务基准，再在固定时间预算下最大化临时任务接纳人数，最后继续优化次级指标。160名临时人员全部接纳，其中阶段B在35114 min预算和当前有限时间列池内获得最优性认证；阶段C最终得到飞机使用时间34033 min、204架次、燃油消耗273629 kg，并启用22架飞机。

三问最终方案均通过独立校验程序复核人员完整性、容量、燃油、路线及相应时间约束；问题一的加油候选数量敏感性实验与问题三的时间列密度实验进一步检验了模型稳定性。结果表明，该框架能够有效协调路线共享、动态容量与机队时空资源，为复杂海上直升机运输计划提供可执行的优化方案。
\end{abstract}

\noindent\textbf{关键词：}海上人员运输；异构直升机；整数规划；词典序优化；机队排班
